\documentclass[10pt,a4paper]{amsart}
\usepackage[margin=19mm]{geometry}
\usepackage{amsmath,amssymb,mathtools}
\usepackage[scr=rsfs,cal=euler]{mathalfa}
\usepackage{xfrac}
\usepackage[unicode,hidelinks,bookmarks=false]{hyperref}
\newcommand{\ie}{i.e.\ }
\newcommand{\cf}{cf.\ }
\newcommand{\resp}{respectively\ }
\newcommand{\Subsection}{\subsection}
\providecommand{\nobreakdash}{\nobreak\hbox{-}}
\setlength{\emergencystretch}{2em}
\multlinegap0pt
\usepackage{amssymb}
\usepackage{dsfont}
\usepackage{mathtools}
\usepackage{xcolor}
\usepackage{tikz}
\usepackage{enumerate}
\newtheorem{theorem}{Theorem}
\DeclareMathOperator{\Aut}{Aut}
\newcommand{\GL}{\text{GL}}
\DeclareMathOperator{\Tors}{Tors}
\newcommand{\f}{\mathfrak{f}}
\newcommand{\mat}[4]{\left(\begin{array}{cc}
#1 &#2\\
#3 &#4
\end{array}\right)}
\newcommand{\p}{\mathfrak{p}}
\title{Torsion points on \(\GL_2\)-type abelian varieties}
\author{Jessica Alessandr\`i, Nirvana Coppola}
\date{Source: arXiv:2602.21047v3 (CC BY 4.0)}
\begin{document}
\maketitle

\noindent\emph{Source and licence.} Torsion points on \(\GL_2\)-type abelian varieties. Version arXiv:2602.21047v3 (CC BY 4.0), distributed by arXiv under the Creative Commons Attribution 4.0 International licence, http://creativecommons.org/licenses/by/4.0/, as recorded in the arXiv licence metadata for this exact version and shown on the abstract page https://arxiv.org/abs/2602.21047v3. Full licence notice: \texttt{licenses/arxiv-2602.21047-CC-BY-4.0.txt}.\par\smallskip

\noindent\textit{Editorial scope.} Counted unit: the article's theorem thm:main\_new (source0.tex lines 260--270). The article's complete original proof of it is delivered with it (source0.tex lines 271--279, one \texttt{proof} environment of 9 lines). No mathematics is written, repaired or re-derived: the statement and the proof are the article's own lines, in the article's own notation, and no retained line is substituted or re-flowed.  Retained uncounted as the source-local context the proof needs: lines 245--256.  Nothing was omitted from any retained span, so the package carries no omission record.  No retained line contains a citation, so the wrapper carries no bibliography and no bibliography entry.  No retained line contains a figure environment, an image-inclusion command or drawing code, so no image is delivered and the package declares no asset dependency.  Editorial formatting only, in addition: an AMS \texttt{amsart} preamble that restates the article's own declarations and macros used by the retained lines, namely the environments \texttt{theorem} on separate counters, together with the standard packages those lines need.  The retained environments print in this document as Theorem 1 in order of appearance, whereas the article prints the same items with the numbering of its own full document.  The complete native source of the article as distributed in the arXiv e-print for this exact version is bundled unchanged as \texttt{sources/195240/source0.tex}.
\begin{theorem}\label{thm:Katz}
    Let \(L\) be a field which is complete under a discrete valuation, and whose residue field is finite. Let \(\mathcal{O}_L\) denote the ring of integers in \(L\), and let \(\pi\) denote a uniformising parameter. Let \(V\) be a 2-dimensional \(L\)-vector space, \(G \subseteq \Aut_L(V)\) a compact subgroup, and \(n \ge 1\) an integer. Suppose that we have the congruence
    \[
    \det(1-g) \equiv 0 \mod \pi^n \quad \forall g \in G.
    \]

    Then there exists a \(L\)-basis \(\{e_1,e_2\}\) of \(V\) and integers \(a,b \ge 0\) with \(a+b=n\) such that the elements of \(G\) expressed as matrices with respect to this basis lie in the subgroup
    \[
    \left\lbrace \mat{1 + \pi^a \mathcal{O}_L}{\pi^a \mathcal{O}_L}{\pi^b \mathcal{O}_L}{1 + \pi^b \mathcal{O}_L} \right\rbrace
    \]
    of \(\GL_2(\mathcal{O}_L)\).
\end{theorem}

\begin{theorem}\label{thm:main_new}
    Let \(A\) be an abelian variety over a number field \(K\) of \(\GL_2\)-type. Let \(\ell\) be a rational prime and let \(\lambda\) be a prime of \(F\) over \(\ell\). Suppose that, for an integer \(n \ge 1\), the congruence
    \[
    P_\p^\lambda(1) \equiv 0 \pmod {\lambda^n}
    \]
    holds for a density one set of finite places \(\p\) of \(K\). Then there exists an abelian variety \(A'\), \(K\)-isogenous to \(A\) through a \(\ell\)-power isogeny, such that 
    \[
    |\Tors(A')(K)| \equiv 0 \pmod {\ell^{\f_{\lambda}n}},
    \]
    where \(\f_\lambda\) is the inertia degree of \(\lambda \mid \ell\).
\end{theorem}

\begin{proof}
    With the notation of Theorem \ref{thm:Katz}, let \(L=F_\lambda\), so that \(\pi=\lambda\), \(V=V_\lambda (A)\) and \(G=\rho_\lambda(G_K)\). The condition \(\det(1-g)\equiv 0 \pmod {\lambda^n}\) for all \(g \in G\) is equivalent to \(P_\p^{\lambda}(1) \equiv 0 \pmod{\lambda^n}\) for a density one set of \(\p\), by \v Cebotarev Density Theorem. Then, by Theorem \ref{thm:Katz}, there exist two sublattices \(\mathcal{L'} \subseteq \mathcal{L}\) of \(V\), such that \(G_K\) acts trivially on the quotient. Such quotient has order \(\ell^{\f_\lambda n}\), indeed in the basis \(\{e_1,e_2\}\) given by the theorem, \(\mathcal{L} \cong \mathcal{O}_{F_\lambda} e_1 + \mathcal{O}_{F_\lambda} e_2\); \(\mathcal{L}' \cong \mathcal{O}_{F_\lambda} \lambda^a e_1 + \mathcal{O}_{F_\lambda} \lambda^b e_2\), with \(a+b=n\), and the quotient is
    \[
    \mathcal{L}/\mathcal{L}' \cong \dfrac{\mathcal{O}_{F_\lambda} e_1 + \mathcal{O}_{F_\lambda} e_2}{\mathcal{O}_{F_\lambda} \lambda^a e_1 + \mathcal{O}_{F_\lambda} \lambda^b e_2} \cong \dfrac{{\mathcal{O}_{F_\lambda}}}{\mathcal{O}_{F_\lambda} \lambda^a} \oplus \dfrac{\mathcal{O}_{F_\lambda}}{\mathcal{O}_{F_\lambda} \lambda^b},
    \]
    which has size \(|\mathbb{F}_\lambda|^{a+b} = \ell^{\f_\lambda n}\).

    Now, the lattice \(\mathcal{L}'\) is \(T_\lambda (A')\) for some abelian variety \(A'\) with \(\lambda\)-adic Galois representation \(\rho_\lambda\). In particular, \(A'\) is \(K\)-isogenous to \(A\) with an isogeny of degree a power of \(\ell\). Therefore, \(\Tors (A')(K)\) contains a subgroup of order \(\ell^{\f_\lambda n}\).
\end{proof}

\end{document}
